% Bewertungspaket zum Gesamttest «Terme und Faktorisieren» — Quelle des PDFs (python3 scripts/build-lp-pdf.py terme-faktorisieren).
% Ganze Punkte je Teilschritt; Werte und Fehlerwerte mit python3 nachgerechnet (scripts/lp/terme-faktorisieren/zahlen.py, 09.10.2026).
% Gerüst aus dem Bewertungspaket Bruchrechnen (nach dessen Prüfung): Auftrag an die KI, Folgefehler, (E)-Zeilen, Restpunkte.
\documentclass[11pt]{article}
\usepackage{../lp-druck}
\renewcommand{\lpname}{Terme und Faktorisieren}
\renewcommand{\lpteilgebiet}{Arithmetik/Algebra ·}
\renewcommand{\lpfuss}{Bewertungspaket}
\newcolumntype{P}{>{\centering\arraybackslash}p{10mm}}
\newenvironment{raster}{\par\smallskip\noindent\renewcommand{\arraystretch}{1.15}\linespread{1.35}\selectfont\tabularx{\linewidth}{@{}>{\raggedright\arraybackslash}p{38mm}P X@{}}\toprule
  \small\textsc{Teilschritt} & \small\textsc{P} & \small\textsc{Musterlösung}\\\midrule}{\bottomrule\endtabularx\par}
\newcommand{\fehler}[1]{\par\smallskip{\small\textcolor{lprot}{\bfseries Typische Fehler:} #1}}
\newcommand{\E}{\,{\footnotesize\bfseries\color{lpgruen}(E)}}
\begin{document}
\lpkopf{Bewertungspaket zum Gesamttest}{}

\begin{lpachtung}{Erst lösen, dann öffnen}
Dieses Paket enthält die Musterlösung. Wer es vor dem Lösen liest, misst am Ende nur, wie gut das Abschreiben gelingt.
\end{lpachtung}

\section*{1 · So gehst du vor}
\begin{enumerate}[leftmargin=*,itemsep=2pt]
\item \textbf{Gesamttest lösen} — vollständig, mit Rechenweg, ohne dieses Paket und ohne Taschenrechner.
\item \textbf{Selbst bewerten} mit Abschnitt~3: Zeile für Zeile vergleichen und die Punkte eintragen. Das geht ohne Hilfsmittel
  und ist der Normalfall.
\item \textbf{Oder bewerten lassen.} Ob du dafür eine KI nutzt und welche, entscheidest du selbst und in eigener Verantwortung —
  je nachdem, was dir aufgrund deines Alters und deiner persönlichen Situation erlaubt ist (Altersgrenzen und
  Nutzungsbedingungen des Anbieters, allenfalls Einverständnis deiner Eltern). Fotografiere oder scanne deine Lösung gut lesbar,
  \textbf{ohne Namen}. Handyfotos tragen oft unsichtbar Standort, Zeit und Gerät mit (Metadaten) — lade lieber einen Scan oder ein
  Bildschirmfoto des Fotos hoch. Lade nur deine Lösung, das Testblatt und dieses PDF hoch, keine persönlichen Daten, und füge den Auftrag aus
  Abschnitt~2 ein.
\item \textbf{Rückmeldung prüfen.} Stimmt, was die KI aus deiner Handschrift gelesen hat? Eine KI kann sich irren —
  im Zweifel gilt das Raster in Abschnitt~3.
\item \textbf{Gezielt wiederholen} nach Abschnitt~4.
\end{enumerate}

\section*{2 · Auftrag an die KI}
\begin{tcolorbox}[colback=lplinie!25,colframe=lplinie,boxrule=0.5pt,arc=1mm,fontupper=\small]
Du bist Korrektorin oder Korrektor für Mathematik auf der Sekundarstufe I in der Schweiz. Im Anhang findest du (1)~das Testblatt mit den Aufgaben, (2)~das Bewertungspaket zum Gesamttest «Terme und Faktorisieren» mit Musterlösung und Punkteraster und (3)~Fotos meiner handschriftlichen Lösung.\par\medskip
Geh so vor:\par
1. Schreib zu jeder Aufgabe G1 bis G9 zuerst ab, was du in meiner Lösung liest — Rechenweg und Ergebnis. Was du nicht sicher lesen kannst, markierst du mit [unsicher]. Nicht raten. Achte besonders auf Minuszeichen, Hochzahlen und Klammern.\par
2. Vergib die Punkte genau nach dem Raster im Bewertungspaket (Abschnitt 3), ganze Punkte je Zeile. Ziehe einen Fehler nur einmal ab: Rechne ich mit einem falschen Zwischenergebnis richtig weiter, gibt es die Folgepunkte. Ausnahme: Zeilen mit \textbf{(E)} gibt es nur für das richtige Ergebnis, nie als Folgepunkt. Die Zeilen «Typische Fehler» sagen, welche Rasterzeile bei diesem Fehler 0 Punkte gibt — sie sind kein zusätzlicher Abzug.\par
3. Andere richtige Lösungswege zählen voll, zum Beispiel Ausmultiplizieren statt binomischer Formel oder eine andere Reihenfolge der Glieder und der Faktoren. Gleichwertig: $(x + 3)^2$ und $(x + 3)(x + 3)$; $-x + 6$ und $6 - x$. Ein Ergebnis zählt nur \textbf{vollständig vereinfacht} (zusammengefasst, ohne Klammern) bzw. beim Faktorisieren \textbf{vollständig faktorisiert}: Ein gleichwertiger, aber nicht fertiger Term gibt in der (E)-Zeile 0 Punkte.\par
4. Begründe jeden Abzug in einem Satz und nenne die Fehlerart (zum Beispiel «Minus nur beim ersten Glied», «mittleres Glied fehlt», «nicht vollständig faktorisiert»). Schreib die Musterlösung nicht ab — zeig mir, wo mein Fehler liegt.\par
5. Gib zum Schluss eine Tabelle «Aufgabe | Punkte | Begründung», die Gesamtpunktzahl (von 25) und — nach der Selbsteinschätzung in Abschnitt 4 — welches Kapitel des Leitprogramms ich wiederholen soll. Keine Note.\par
6. Fehlt ein Foto oder ist eine Aufgabe unlesbar, sag es, statt sie zu bewerten.
\end{tcolorbox}

\needspace{16\baselineskip}
\section*{3 · Musterlösung und Punkteraster}
\textbf{Allgemein:} Ganze Punkte je Zeile. Ein Fehler wird nur einmal abgezogen (Folgefehler). Gleichwertige Wege zählen voll.
In der Spalte P: \textbf{(E)}~= Ergebnispunkt, Ergebnis richtig, vollständig vereinfacht bzw.\ vollständig faktorisiert \emph{und}
Weg erkennbar; ohne Zeichen = die Zeile bewertet den Weg oder die Begründung. Der ganze Test ist ohne Taschenrechner gestellt.
«Typische Fehler» nennt, welche Zeile 0 Punkte gibt, und die Punkte, die bleiben.
\textbf{Teilrichtige Zeilen:} Verlangt eine Zeile mehrere Angaben und stimmt nur ein Teil, gibt die Zeile 0 Punkte,
ausser die Zeile nennt ausdrücklich eine Aufteilung.
\textbf{Schreibweise:} Die Reihenfolge der Glieder und Faktoren ist frei ($-x + 6 = 6 - x$; $(x + 4)(x - 6) = (x - 6)(x + 4)$).
Mal-Punkte dürfen fehlen. Beim Faktorisieren zählt nur ein \emph{Produkt}; eine gleichwertige Summe gibt den Ergebnispunkt nicht.

\begin{lpaufgabe}{G1}{Klammern auflösen}{3 P}
\begin{raster}
Klammer mit Zahl & 1 & $-2(3x - 4) = -6x + 8$: beide Glieder multipliziert, Vorzeichen richtig.\\
Zweite Klammer & 1 & $+(x - 1) = x - 1$, und die $x$-Glieder zusammengefasst: $5x - 6x + x = 0$.\\
Ergebnis & 1\E & $5x - 6x + 8 + x - 1 = 7$. Auch $0x + 7$ zählt, wenn am Schluss $7$ steht.\\
\end{raster}
\fehler{$-2(3x - 4) = -6x - 8$ (Vorzeichen), Ergebnis $-9$: erste Zeile und Ergebniszeile 0 $\to$ 1 von 3\,P.
$-2(3x - 4) = -6x - 4$ (nur das erste Glied multipliziert), Ergebnis $-5$: erste Zeile und Ergebniszeile 0 $\to$ 1 von 3\,P.
Richtig bis $5x - 6x + 8 + x - 1$, dann verrechnet: Ergebniszeile 0 $\to$ 2 von 3\,P.}
\end{lpaufgabe}

\begin{lpaufgabe}{G2}{Minus vor der Klammer}{2 P}
\begin{raster}
Klammer & 1 & $3a^2 - a^2 + 2a - 2a$: Das Minus dreht \emph{beide} Vorzeichen um.\\
Ergebnis & 1\E & $2a^2$\\
\end{raster}
\fehler{Minus nur beim ersten Glied ($3a^2 - a^2 - 2a - 2a = 2a^2 - 4a$): beide Zeilen 0 $\to$ 0 von 2\,P.
Ungleichartiges zusammengefasst ($2a^2 + 2a - 2a = 2a^4$ o.\,ä.): Ergebniszeile 0 $\to$ 1 von 2\,P.}
\end{lpaufgabe}

\begin{lpaufgabe}{G3}{Klammer mal Klammer}{3 P}
\begin{raster}
Äussere Produkte & 1 & $2x \cdot x = 2x^2$ und $5 \cdot (-3) = -15$\\
Gemischte Produkte & 1 & $2x \cdot (-3) = -6x$ und $5 \cdot x = 5x$\\
Ergebnis & 1\E & $2x^2 - 6x + 5x - 15 = 2x^2 - x - 15$\\
\end{raster}
\fehler{Nur vorne mal vorne, hinten mal hinten ($2x^2 - 15$): Zeile «gemischte Produkte» und Ergebniszeile 0 $\to$ 1 von 3\,P.
Vorzeichenfehler bei einem gemischten Produkt (z.\,B. $2x^2 + 11x - 15$): Zeile «gemischte Produkte» und Ergebniszeile 0 $\to$ 1 von 3\,P.
$+15$ statt $-15$ am Schluss: Zeile «äussere Produkte» und Ergebniszeile 0 $\to$ 1 von 3\,P.
Vier Produkte richtig, nicht zusammengefasst ($2x^2 - 6x + 5x - 15$ als Schluss): Ergebniszeile 0 $\to$ 2 von 3\,P.}
\end{lpaufgabe}

\begin{lpaufgabe}{G4}{Binomische Formel}{3 P}
\begin{raster}
Quadrate & 1 & $(3y)^2 = 9y^2$ und $4^2 = 16$\\
Mittleres Glied & 1 & $-2 \cdot 3y \cdot 4 = -24y$\\
Ergebnis & 1\E & $9y^2 - 24y + 16$. Ausmultipliziert als $(3y - 4)(3y - 4)$ zählt gleich.\\
\end{raster}
\fehler{Mittleres Glied vergessen ($9y^2 + 16$ oder $9y^2 - 16$): Zeile «mittleres Glied» und Ergebniszeile 0 $\to$ 1 von 3\,P.
$3y^2$ statt $9y^2$ (Zahl vor $y$ nicht quadriert): Zeile «Quadrate» und Ergebniszeile 0 $\to$ 1 von 3\,P.
$-12y$ statt $-24y$ (Faktor 2 vergessen) oder $+24y$ (Vorzeichen): Zeile «mittleres Glied» und Ergebniszeile 0 $\to$ 1 von 3\,P.}
\end{lpaufgabe}

\begin{lpaufgabe}{G5}{Kopfrechnen}{2 P}
\begin{raster}
Formel & 1 & $31 \cdot 29 = (30 + 1)(30 - 1)$, dritte binomische Formel $(a + b)(a - b) = a^2 - b^2$\\
Ergebnis & 1\E & $30^2 - 1^2 = 900 - 1 = 899$\\
\end{raster}
\fehler{Schriftlich multipliziert, ohne Formel: Formelzeile 0 $\to$ 1 von 2\,P.
$900 + 1 = 901$ oder $(30 + 1)^2$ gerechnet: Ergebniszeile 0; die Zerlegung $(30 + 1)(30 - 1)$ gibt die Formelzeile.}
\end{lpaufgabe}

\begin{lpaufgabe}{G6}{Ausklammern}{3 P}
\begin{raster}
Gemeinsamer Faktor & 1 & $6a$: $\text{ggT}(12;\ 18) = 6$, und beide Glieder enthalten $a$.\\
Klammerinhalt & 1 & $12a^2 : 6a = 2a$ und $18a : 6a = 3$ (zum gewählten Faktor passend; Probe stimmt)\\
Ergebnis & 1\E & $6a(2a + 3)$\\
\end{raster}
\fehler{Nicht alles ausgeklammert ($6(2a^2 + 3a)$, $a(12a + 18)$, $3a(4a + 6)$, $2a(6a + 9)$): Faktorzeile und Ergebniszeile 0;
der Klammerinhalt passt zum gewählten Faktor $\to$ 1 von 3\,P.
Klammerinhalt falsch ($6a(2a + 18)$): Inhalts- und Ergebniszeile 0 $\to$ 1 von 3\,P.}
\end{lpaufgabe}

\begin{lpaufgabe}{G7}{Zweiklammersatz}{3 P}
\begin{raster}
Ansatz & 1 & zwei Zahlen mit Produkt $-24$ und Summe $-2$ gesucht (Paare, Tabelle oder Probieren mit Probe)\\
Zahlen & 1 & $-6$ und $4$: $(-6) \cdot 4 = -24$, $-6 + 4 = -2$\\
Ergebnis & 1\E & $(x - 6)(x + 4)$\\
\end{raster}
\fehler{Vorzeichen vertauscht ($(x + 6)(x - 4) = x^2 + 2x - 24$): Zahlen- und Ergebniszeile 0 $\to$ 1 von 3\,P.
Summe statt Produkt verwechselt (z.\,B. $(x - 8)(x + 3) = x^2 - 5x - 24$): Zahlen- und Ergebniszeile 0 $\to$ 1 von 3\,P.
Gleichung gelöst statt faktorisiert («$x = 6$ oder $x = -4$»): Ergebniszeile 0; die Zahlen $-6$ und $4$ sind erkennbar $\to$ 2 von 3\,P.}
\end{lpaufgabe}

\begin{lpaufgabe}{G8}{Vollständig faktorisieren}{3 P}
\begin{raster}
Ausklammern & 1 & $3x^2 - 12 = 3(x^2 - 4)$\\
Dritte Formel & 1 & $x^2 - 4 = (x - 2)(x + 2)$\\
Ergebnis & 1\E & $3(x - 2)(x + 2)$. Gleichwertig: $3(x + 2)(x - 2)$.\\
\end{raster}
\fehler{Nur $3(x^2 - 4)$: Formel- und Ergebniszeile 0 $\to$ 1 von 3\,P.
$(3x - 6)(x + 2)$ oder $(x - 2)(3x + 6)$ (die $3$ nicht ausgeklammert, nicht vollständig): Ausklammerzeile und Ergebniszeile 0 $\to$ 1 von 3\,P.
$3(x - 2)^2$ (Binom verwechselt): Formel- und Ergebniszeile 0 $\to$ 1 von 3\,P.}
\end{lpaufgabe}

\begin{lpaufgabe}{G9}{Fehler finden}{3 P}
\begin{raster}
Begründung & 1 & Zuletzt wird $+\,9$ gerechnet: Der ganze Term ist eine \emph{Summe}, kein Produkt. Nur «falsch» ohne Grund: 0.\\
Faktorisierung & 1\E & $x^2 + 6x + 9 = (x + 3)^2$, denn $9 = 3^2$ und $6x = 2 \cdot x \cdot 3$\\
Probe & 1 & $(x + 3)^2 = x^2 + 6x + 9$ ausmultipliziert. Die Probe zählt auch als Folgepunkt, wenn sie einen eigenen falschen Vorschlag richtig ausmultipliziert und so den Fehler zeigt.\\
\end{raster}
\fehler{$(x + 3)(x - 3)$ (Binom verwechselt): Faktorisierungszeile 0; eine richtige Probe ($x^2 - 9$, also nicht gleich) gibt ihren Punkt $\to$ höchstens 2 von 3\,P.
$(x + 9)(x + 1)$ o.\,ä.: Faktorisierungszeile 0.}
\end{lpaufgabe}

\needspace{12\baselineskip}
\section*{4 · Selbsteinschätzung}
\begin{tabularx}{\linewidth}{@{}l X@{}}
\toprule
\textbf{22 – 25 P} & Die geprüften Teile sitzen. Wo du Punkte verloren hast: das Kapitel dieser Aufgabe nochmals.\\
\textbf{17 – 21 P} & Den schwächsten Teil nochmals, gemessen am Anteil (Teil A 5, B 8, C 12 P): Tüfteln und Übungen der Kapitel dieses Teils.\\
\textbf{11 – 16 P} & Zurück zu den Kapiteln aller Aufgaben, in denen du Punkte verloren hast.\\
\textbf{0 – 10 P} & Zurück zu den Kapiteln jedes Teils, in dem du weniger als die Hälfte erreicht hast; sind es alle drei, zum Vortest und zu Kapitel 1.\\
\bottomrule
\end{tabularx}

\smallskip
\textbf{Aufgabe $\to$ Kapitel:} G1, G2 $\to$ 1; G3 $\to$ 2; G4, G5 $\to$ 3; G6 bis G9 $\to$ 4.
\end{document}
